% ============================================================
%  CAPÍTULO 1 — INTRODUCCIÓN
% ============================================================
\chapter{Introducción}

En Colombia, las \ac{MAP} han sido responsables de miles de accidentes,
generando miedo en las comunidades principalmente rurales y frenando su
desarrollo productivo y social. Colombia es considerado uno de los países más
afectados en el mundo por esta problemática: 715 municipios han sido
contaminados con minas, aunque 523 de ellos ya han sido declarados libres
\cite{reuters2025}. Estos dispositivos explosivos, utilizados por grupos armados
ilegales durante décadas de conflicto, han dejado más de 12\,673 víctimas entre muertos y heridos desde 1990~\cite{aicma2026}
. Además, la persistencia de estas
minas en el país ha conllevado el desplazamiento y confinamiento de la población
civil; según datos del 2023, más de 14\,700 personas tuvieron que abandonar su
hogar temporalmente o limitar su movilidad debido a la presencia de \ac{MAP} u
otros artefactos explosivos en su territorio \cite{eltiempo2024}. Pese a los
avances en el posconflicto y la reducción anual de accidentes, el peligro sigue
presente; en 2025 se registraron 136 accidentes, un incremento del 23\,\% 
respecto al año anterior~\cite{aicma2026}.

Las entidades encargadas del desminado humanitario en Colombia son
principalmente las Fuerzas Militares (Ejército Nacional) y las \ac{OCDH},
las cuales se rigen por la normativa técnica de la \ac{AICMA}, que usan
tres metodologías para la eliminación de artefactos explosivos: la \ac{TDM},
que emplea detectores de metales y herramientas especializadas; la \ac{TDC};
y la \ac{TDEM}, destinada a la destrucción controlada de los explosivos
localizados \cite{forero2022}. Los métodos de detección más usados incluyen
los \ac{MD}, los \ac{GPR} --- que emiten pulsos electromagnéticos para detectar
objetos enterrados --- y el uso de animales entrenados. Si bien estas técnicas
han sido efectivas en muchos casos, presentan limitaciones importantes: los
\ac{MD} son poco eficaces con minas artesanales de contenido metálico mínimo;
los \ac{GPR} requieren operar muy cerca del terreno peligroso; y el uso de
animales los expone a posibles detonaciones.

Frente a este panorama, los trabajos previos del grupo de investigación \ac{PSI}
de la Universidad del Valle \cite{forero2022,tenorio2024} demostraron que las
imágenes termográficas adquiridas desde un \ac{UAV} a baja altura permiten
detectar \ac{MAP} tipo quiebrapatas con porcentajes de éxito superiores al 88\,\%
en condiciones controladas. El enfoque térmico, sin embargo, depende de diferencias transitorias de
temperatura entre la mina y el suelo. Esas diferencias se ven afectadas
por la hora del día, la humedad del suelo y la radiación solar.

Las imágenes multiespectrales ofrecen una alternativa complementaria: el
enterramiento de una \ac{MAP} perturba el suelo superficial e induce estrés
en la vegetación circundante, generando firmas espectrales anómalas más
persistentes que la señal térmica \cite{barnawi2022}. Frente a las imágenes
hiperespectrales, las multiespectrales son más accesibles en costo y
procesamiento, ya que capturan un número reducido de bandas orientadas a
discriminar materiales de interés \cite{makki2017}.

Los resultados de este trabajo respaldan esa estabilidad: el sistema
alcanza \ac{AUC} superiores a 0.80 en sesiones capturadas en distintas
franjas horarias ---mañana temprana, mediodía y atardecer---. El enfoque
térmico, en cambio, requiere condiciones específicas de irradiancia y
temperatura para generar el contraste detectable \cite{tenorio2024}.



Por lo anterior, el presente trabajo propone el desarrollo de un sistema de visión
artificial capaz de identificar \ac{MAP} enterradas superficialmente en suelo
arenoso con vegetación limitada, a partir de imágenes multiespectrales adquiridas
a \SI{1}{m} de altura con una cámara MicaSense RedEdge-MX de cinco bandas
espectrales. El experimento se realizó en condiciones controladas: las minas
emuladas (réplicas sin material explosivo) se enterraron a cuatro profundidades (1, 3, 5 y 7~cm) y en cada zona
se dispusieron objetos de interferencia \textbf{enterrados a la misma profundidad}
(roca, botella de vidrio y lata de aluminio) para dotar a la base de datos de
mayor realismo de cara a trabajos futuros. El sistema implementado aborda la
clasificación binaria \textit{mina} / \textit{no mina}.

\section{Objetivos}

\subsection*{Objetivo general}

Desarrollar un sistema de visión artificial, capaz de identificar \ac{MAP}
enterradas superficialmente en suelo arenoso con vegetación limitada, a partir de
imágenes multiespectrales adquiridas a \SI{1}{m} de altura.

\subsection*{Objetivos específicos}

\begin{enumerate}
  \item Construir una base de datos de imágenes multiespectrales de minas
    enterradas en suelo real y vegetación limitada a \SI{1}{m} de altura.
  \item Implementar técnicas de procesamiento en imágenes multiespectrales para
    mejorar la discriminación de las \ac{MAP} respecto al resto de componentes
    en las imágenes.
  \item Implementar una técnica de reconocimiento de patrones para identificar la
    presencia de una \ac{MAP}.
  \item Establecer los alcances y limitaciones del sistema desarrollado.
\end{enumerate}

\section{Justificación}

La realización de este trabajo está justificada por su importancia humanitaria,
social y científica. En términos humanitarios, aportará al esfuerzo por reducir
víctimas de \ac{MAP}, un problema que ha cobrado más de 12\,673 víctimas en Colombia desde 1990~\cite{aicma2026}
. Una tecnología de detección
remota eficaz contribuiría a salvar vidas —tanto de civiles como de operarios—
al minimizar la necesidad de que el personal de desminado ingrese directamente
a las áreas minadas. Asimismo, facilitar el desminado agilizará la restitución de
tierras seguras para las comunidades rurales, permitiendo reactivar proyectos
agrícolas y vías de comunicación en zonas que hoy permanecen improductivas
por el temor a las minas. En línea con los compromisos del Estado colombiano y
tratados internacionales como la Convención de Ottawa, esta iniciativa tecnológica
sumará esfuerzos hacia la meta de un territorio libre de minas, en concordancia
con diez de los diecisiete \ac{ODS} \cite{jodi2021}.

Desde el punto de vista científico, Colombia registra 51 publicaciones sobre
detección de \ac{MAP} en Scopus, pero ninguna de ellas emplea imágenes
multiespectrales como técnica de detección \cite{IsaacAnteproyecto}. Este trabajo
abre esa línea de investigación en el contexto colombiano, complementando los
avances previos del grupo \ac{PSI} basados en termografía.

\section{Estructura del documento}

El Capítulo~2 revisa el estado del arte en detección de \ac{MAP} y establece
las bases teóricas de las imágenes multiespectrales y los algoritmos de
clasificación empleados. El Capítulo~3 describe el diseño experimental y el
\textit{pipeline} de procesamiento, desde la captura hasta la extracción de
características. Los resultados se presentan en el Capítulo~4, y el
Capítulo~5 recoge las conclusiones, los alcances y limitaciones del sistema,
y las líneas de trabajo futuro.
